% 决赛展示 PPT 公共样式（Beamer，16:9，XeLaTeX）——商业路演稿标准
% 用法：
%   \documentclass[aspectratio=169,10pt,t]{beamer}   % t：内容顶端对齐，结论条压到页底
%   % 决赛展示 PPT 公共样式（Beamer，16:9，XeLaTeX）——商业路演稿标准
% 用法：
%   \documentclass[aspectratio=169,10pt,t]{beamer}   % t：内容顶端对齐，结论条压到页底
%   % 决赛展示 PPT 公共样式（Beamer，16:9，XeLaTeX）——商业路演稿标准
% 用法：
%   \documentclass[aspectratio=169,10pt,t]{beamer}   % t：内容顶端对齐，结论条压到页底
%   % 决赛展示 PPT 公共样式（Beamer，16:9，XeLaTeX）——商业路演稿标准
% 用法：
%   \documentclass[aspectratio=169,10pt,t]{beamer}   % t：内容顶端对齐，结论条压到页底
%   \input{../_style/slides_preamble.tex}
%   \title{…} \subtitle{…} \author{fx战队} \date{2026 年 11 月}
%   \TrackName{因子赛道}
%
% 页面结构：结论标题（浅显） + 主体（图 / 示意图 / 表 / 要点块） + 底部结论条 \takeaway{…}
% 组件：\ChapterFrame、\keynum、\point、panel 环境、\takeaway、\hl

\usepackage[fontset=none,scheme=plain]{ctex}
% 字体：微软雅黑（取自本机 Office 安装包，不复制进仓库），大号数字用 DIN
\def\yhpath{/Applications/Microsoft PowerPoint.app/Contents/Resources/DFonts/}
\setCJKmainfont{msyh.ttc}[Path=\yhpath,BoldFont=msyhbd.ttc]
\setCJKsansfont{msyh.ttc}[Path=\yhpath,BoldFont=msyhbd.ttc]
\setCJKmonofont{msyh.ttc}[Path=\yhpath]
\newCJKfontfamily\titlecjk{msyhbd.ttc}[Path=\yhpath]
\setsansfont{msyh.ttc}[Path=\yhpath,BoldFont=msyhbd.ttc]
\setmainfont{msyh.ttc}[Path=\yhpath,BoldFont=msyhbd.ttc]
\newfontfamily\titlelatin{msyhbd.ttc}[Path=\yhpath]
\newfontfamily\numfont{DIN Alternate}
\usefonttheme{professionalfonts}

\usepackage{newunicodechar}
\newunicodechar{→}{\ensuremath{\rightarrow}}
\newunicodechar{↑}{\ensuremath{\uparrow}}
\newunicodechar{↓}{\ensuremath{\downarrow}}
\newunicodechar{≈}{\ensuremath{\approx}}
\newunicodechar{≤}{\ensuremath{\leq}}
\newunicodechar{≥}{\ensuremath{\geq}}
\newunicodechar{−}{\ensuremath{-}}
\newunicodechar{×}{\ensuremath{\times}}
\newunicodechar{÷}{\ensuremath{\div}}
\newunicodechar{±}{\ensuremath{\pm}}
\newunicodechar{β}{\ensuremath{\beta}}
\newunicodechar{λ}{\ensuremath{\lambda}}

% 配色：深蓝主色、橙色强调、浅蓝底色
\definecolor{snavy}{HTML}{0F2E5A}
\definecolor{sblue}{HTML}{2F6DB5}
\definecolor{stint}{HTML}{E8EFF8}
\definecolor{sorange}{HTML}{E8772E}
\definecolor{sgrey}{HTML}{5E6672}
\definecolor{sline}{HTML}{C9D3E0}
\definecolor{sink}{HTML}{1E2329}

\setbeamercolor{normal text}{fg=sink,bg=white}
\setbeamercolor{structure}{fg=snavy}
\setbeamercolor{frametitle}{fg=snavy}
\setbeamercolor{alerted text}{fg=sorange}
\setbeamertemplate{navigation symbols}{}
\setbeamertemplate{itemize items}{\textcolor{sblue}{\rule[0.25ex]{0.42em}{0.42em}}}
\setlength{\leftmargini}{1.1em}
\linespread{1.18}
\setbeamersize{text margin left=1.0cm, text margin right=1.0cm}
% 内容顶端对齐，结论条用 \vfill 压到页底
\makeatletter\beamer@centeredfalse\makeatother

\providecommand{\trackname}{}
\gdef\chaplabel{}
\newcommand{\TrackName}[1]{\renewcommand{\trackname}{#1}}

\usepackage{tikz}
\usetikzlibrary{arrows.meta,positioning,calc,fit,decorations.pathreplacing,shapes.geometric}

% 顶部深蓝细带 + 标题左侧橙色色块 + 右上章节标签
\setbeamertemplate{frametitle}{%
  \begin{tikzpicture}[remember picture,overlay]
    \fill[snavy] (current page.north west) rectangle ([yshift=-0.12cm]current page.north east);
  \end{tikzpicture}%
  \vspace*{0.36cm}%
  \begin{minipage}[c]{0.80\textwidth}%
    {\color{sorange}\rule[-0.05cm]{0.13cm}{0.52cm}}\hspace{0.22cm}%
    {\titlelatin\titlecjk\color{snavy}\fontsize{16pt}{20pt}\selectfont\insertframetitle}%
  \end{minipage}\hfill
  \ifx\chaplabel\empty\else
  \begin{tikzpicture}[baseline=-0.6ex]
    \node[fill=snavy,text=white,rounded corners=1.5pt,inner xsep=6pt,inner ysep=3pt,font=\footnotesize]{\chaplabel};
  \end{tikzpicture}\fi\par
  \vspace{0.25cm}}

% 页脚：只有页码
\setbeamertemplate{footline}{\hfill{\numfont\color{sgrey}\small\insertframenumber}\hspace*{0.6cm}\vspace*{0.3cm}}

\usepackage{appendixnumberbeamer}

% ---------- 组件 ----------
\newcommand{\hl}[1]{\textcolor{sorange}{\bfseries #1}}
% 大号关键数字：\keynum{0.290}{换手率（原 0.518）}
\newcommand{\keynum}[2]{%
  \begin{minipage}[t]{\linewidth}\raggedright
    {\numfont\color{sorange}\fontsize{26pt}{28pt}\selectfont #1}\par\vspace{0.06cm}
    {\color{sink}\small #2}\par
  \end{minipage}}
% 编号要点：\point{1}{加粗短句}{一两句解释}
\newcommand{\point}[3]{%
  \begin{tikzpicture}[baseline=(n.base)]
    \node[circle,fill=snavy,text=white,inner sep=1.6pt,minimum size=0.48cm,font=\numfont\small](n){#1};
  \end{tikzpicture}\hspace{0.18cm}%
  \begin{minipage}[t]{\dimexpr\linewidth-0.75cm\relax}
    {\bfseries\color{snavy}#2}\par\vspace{0.03cm}{\small\color{sink}#3}
  \end{minipage}\par\vspace{0.16cm}}
% 浅蓝底面板：\begin{panel}{标题} … \end{panel}
\newenvironment{panel}[1]{%
  \par\noindent\begin{tikzpicture}\node[fill=stint,rounded corners=2pt,inner sep=8pt,text width=\dimexpr\linewidth-16pt\relax,align=flush left]\bgroup
  {\bfseries\color{snavy}#1}\par\vspace{0.12cm}\small}
  {\egroup;\end{tikzpicture}\par}
% 底部结论条：一句话说明这一页证明了什么（固定在页底，正文不要压到这一区域）
\newcommand{\takeaway}[1]{%
  \begin{tikzpicture}[remember picture,overlay]
    \node[fill=snavy,text=white,inner xsep=10pt,inner ysep=6pt,anchor=south,
          text width=\dimexpr\paperwidth-2cm-20pt\relax,align=flush left,font=\small]
      at ([yshift=0.75cm]current page.south) {#1};
  \end{tikzpicture}}

% 表格
\usepackage{booktabs,tabularx,array,colortbl}
\renewcommand{\arraystretch}{1.32}
\newcolumntype{L}{>{\raggedright\arraybackslash}X}
\newcolumntype{C}{>{\centering\arraybackslash}X}
\newcommand{\headrow}{\rowcolor{snavy}}
\newcommand{\thc}[1]{\textcolor{white}{\bfseries #1}}
\newcommand{\ourrow}{\rowcolor{stint}}

% 图形
\usepackage{graphicx}
\usepackage{pgfplots}
\pgfplotsset{compat=1.18,tick label style={font=\numfont\small},label style={font=\small},/pgf/number format/assume math mode=true}
\tikzset{
  node/.style  = {draw=snavy, line width=0.8pt, fill=white, rounded corners=2pt, align=flush center,
                  inner xsep=7pt, inner ysep=5pt, font=\small},
  fillnode/.style = {fill=snavy, text=white, rounded corners=2pt, align=flush center,
                  inner xsep=7pt, inner ysep=5pt, font=\small\bfseries},
  tintnode/.style = {fill=stint, draw=none, rounded corners=2pt, align=flush center,
                  inner xsep=7pt, inner ysep=5pt, font=\small},
  key/.style   = {node, draw=sorange, line width=1.1pt},
  flow/.style  = {-{Stealth[length=2.6mm]}, line width=0.9pt, draw=snavy},
  note/.style  = {font=\small, text=sgrey, align=flush center},
}

% 封面：深蓝满版
\setbeamertemplate{title page}{%
  \begin{tikzpicture}[remember picture,overlay]
    \fill[snavy] (current page.south west) rectangle (current page.north east);
    \fill[sorange] ([xshift=1.2cm,yshift=-2.3cm]current page.north west) rectangle ++(1.6cm,0.12cm);
    \node[anchor=north west,text=white,text width=0.78\paperwidth,align=flush left]
      at ([xshift=1.2cm,yshift=-2.7cm]current page.north west)
      {{\titlecjk\fontsize{25pt}{33pt}\selectfont\inserttitle}\par\vspace{0.45cm}
       {\large\color{white!80!snavy}\insertsubtitle}};
    \node[anchor=south west,text=white,text width=0.6\paperwidth,align=flush left]
      at ([xshift=1.2cm,yshift=1.1cm]current page.south west)
      {{\large\bfseries\insertauthor}\par\vspace{0.12cm}{\small\color{white!75!snavy}\trackname\quad\insertdate}};
  \end{tikzpicture}}

% 章节页：深蓝满版，左侧大号编号与本章要回答的问题，右侧故事进度
\providecommand{\ChapterList}[1]{}
\newcommand{\ChapterFrame}[4]{%
  \section[#1\enspace #2]{#3}\gdef\chaplabel{#1\enspace #2}
  {\setbeamertemplate{footline}{}
  \begin{frame}[plain,noframenumbering]
    \begin{tikzpicture}[remember picture,overlay]
      \fill[snavy] (current page.south west) rectangle (current page.north east);
      \node[anchor=north west,text=white,text width=0.46\paperwidth,align=flush left]
        at ([xshift=1.2cm,yshift=-2.0cm]current page.north west)
        {{\numfont\color{sorange}\fontsize{48pt}{50pt}\selectfont #1}\par\vspace{0.35cm}
         {\titlecjk\fontsize{20pt}{26pt}\selectfont #3}\par\vspace{0.4cm}
         {\normalsize\color{white!85!snavy}#4}};
      \node[anchor=north west,text=white,text width=0.36\paperwidth,align=flush left]
        at ([xshift=0.58\paperwidth,yshift=-2.1cm]current page.north west)
        {\ChapterList{#1}};
    \end{tikzpicture}
  \end{frame}}}
\newcommand{\chapitem}[3]{\ifnum#1=#2{\color{white}\bfseries #2\enspace #3}\else{\color{white!55!snavy}#2\enspace #3}\fi\par\vspace{0.3cm}}

\hypersetup{colorlinks=true, linkcolor=snavy, urlcolor=snavy}

%   \title{…} \subtitle{…} \author{fx战队} \date{2026 年 11 月}
%   \TrackName{因子赛道}
%
% 页面结构：结论标题（浅显） + 主体（图 / 示意图 / 表 / 要点块） + 底部结论条 \takeaway{…}
% 组件：\ChapterFrame、\keynum、\point、panel 环境、\takeaway、\hl

\usepackage[fontset=none,scheme=plain]{ctex}
% 字体：微软雅黑（取自本机 Office 安装包，不复制进仓库），大号数字用 DIN
\def\yhpath{/Applications/Microsoft PowerPoint.app/Contents/Resources/DFonts/}
\setCJKmainfont{msyh.ttc}[Path=\yhpath,BoldFont=msyhbd.ttc]
\setCJKsansfont{msyh.ttc}[Path=\yhpath,BoldFont=msyhbd.ttc]
\setCJKmonofont{msyh.ttc}[Path=\yhpath]
\newCJKfontfamily\titlecjk{msyhbd.ttc}[Path=\yhpath]
\setsansfont{msyh.ttc}[Path=\yhpath,BoldFont=msyhbd.ttc]
\setmainfont{msyh.ttc}[Path=\yhpath,BoldFont=msyhbd.ttc]
\newfontfamily\titlelatin{msyhbd.ttc}[Path=\yhpath]
\newfontfamily\numfont{DIN Alternate}
\usefonttheme{professionalfonts}

\usepackage{newunicodechar}
\newunicodechar{→}{\ensuremath{\rightarrow}}
\newunicodechar{↑}{\ensuremath{\uparrow}}
\newunicodechar{↓}{\ensuremath{\downarrow}}
\newunicodechar{≈}{\ensuremath{\approx}}
\newunicodechar{≤}{\ensuremath{\leq}}
\newunicodechar{≥}{\ensuremath{\geq}}
\newunicodechar{−}{\ensuremath{-}}
\newunicodechar{×}{\ensuremath{\times}}
\newunicodechar{÷}{\ensuremath{\div}}
\newunicodechar{±}{\ensuremath{\pm}}
\newunicodechar{β}{\ensuremath{\beta}}
\newunicodechar{λ}{\ensuremath{\lambda}}

% 配色：深蓝主色、橙色强调、浅蓝底色
\definecolor{snavy}{HTML}{0F2E5A}
\definecolor{sblue}{HTML}{2F6DB5}
\definecolor{stint}{HTML}{E8EFF8}
\definecolor{sorange}{HTML}{E8772E}
\definecolor{sgrey}{HTML}{5E6672}
\definecolor{sline}{HTML}{C9D3E0}
\definecolor{sink}{HTML}{1E2329}

\setbeamercolor{normal text}{fg=sink,bg=white}
\setbeamercolor{structure}{fg=snavy}
\setbeamercolor{frametitle}{fg=snavy}
\setbeamercolor{alerted text}{fg=sorange}
\setbeamertemplate{navigation symbols}{}
\setbeamertemplate{itemize items}{\textcolor{sblue}{\rule[0.25ex]{0.42em}{0.42em}}}
\setlength{\leftmargini}{1.1em}
\linespread{1.18}
\setbeamersize{text margin left=1.0cm, text margin right=1.0cm}
% 内容顶端对齐，结论条用 \vfill 压到页底
\makeatletter\beamer@centeredfalse\makeatother

\providecommand{\trackname}{}
\gdef\chaplabel{}
\newcommand{\TrackName}[1]{\renewcommand{\trackname}{#1}}

\usepackage{tikz}
\usetikzlibrary{arrows.meta,positioning,calc,fit,decorations.pathreplacing,shapes.geometric}

% 顶部深蓝细带 + 标题左侧橙色色块 + 右上章节标签
\setbeamertemplate{frametitle}{%
  \begin{tikzpicture}[remember picture,overlay]
    \fill[snavy] (current page.north west) rectangle ([yshift=-0.12cm]current page.north east);
  \end{tikzpicture}%
  \vspace*{0.36cm}%
  \begin{minipage}[c]{0.80\textwidth}%
    {\color{sorange}\rule[-0.05cm]{0.13cm}{0.52cm}}\hspace{0.22cm}%
    {\titlelatin\titlecjk\color{snavy}\fontsize{16pt}{20pt}\selectfont\insertframetitle}%
  \end{minipage}\hfill
  \ifx\chaplabel\empty\else
  \begin{tikzpicture}[baseline=-0.6ex]
    \node[fill=snavy,text=white,rounded corners=1.5pt,inner xsep=6pt,inner ysep=3pt,font=\footnotesize]{\chaplabel};
  \end{tikzpicture}\fi\par
  \vspace{0.25cm}}

% 页脚：只有页码
\setbeamertemplate{footline}{\hfill{\numfont\color{sgrey}\small\insertframenumber}\hspace*{0.6cm}\vspace*{0.3cm}}

\usepackage{appendixnumberbeamer}

% ---------- 组件 ----------
\newcommand{\hl}[1]{\textcolor{sorange}{\bfseries #1}}
% 大号关键数字：\keynum{0.290}{换手率（原 0.518）}
\newcommand{\keynum}[2]{%
  \begin{minipage}[t]{\linewidth}\raggedright
    {\numfont\color{sorange}\fontsize{26pt}{28pt}\selectfont #1}\par\vspace{0.06cm}
    {\color{sink}\small #2}\par
  \end{minipage}}
% 编号要点：\point{1}{加粗短句}{一两句解释}
\newcommand{\point}[3]{%
  \begin{tikzpicture}[baseline=(n.base)]
    \node[circle,fill=snavy,text=white,inner sep=1.6pt,minimum size=0.48cm,font=\numfont\small](n){#1};
  \end{tikzpicture}\hspace{0.18cm}%
  \begin{minipage}[t]{\dimexpr\linewidth-0.75cm\relax}
    {\bfseries\color{snavy}#2}\par\vspace{0.03cm}{\small\color{sink}#3}
  \end{minipage}\par\vspace{0.16cm}}
% 浅蓝底面板：\begin{panel}{标题} … \end{panel}
\newenvironment{panel}[1]{%
  \par\noindent\begin{tikzpicture}\node[fill=stint,rounded corners=2pt,inner sep=8pt,text width=\dimexpr\linewidth-16pt\relax,align=flush left]\bgroup
  {\bfseries\color{snavy}#1}\par\vspace{0.12cm}\small}
  {\egroup;\end{tikzpicture}\par}
% 底部结论条：一句话说明这一页证明了什么（固定在页底，正文不要压到这一区域）
\newcommand{\takeaway}[1]{%
  \begin{tikzpicture}[remember picture,overlay]
    \node[fill=snavy,text=white,inner xsep=10pt,inner ysep=6pt,anchor=south,
          text width=\dimexpr\paperwidth-2cm-20pt\relax,align=flush left,font=\small]
      at ([yshift=0.75cm]current page.south) {#1};
  \end{tikzpicture}}

% 表格
\usepackage{booktabs,tabularx,array,colortbl}
\renewcommand{\arraystretch}{1.32}
\newcolumntype{L}{>{\raggedright\arraybackslash}X}
\newcolumntype{C}{>{\centering\arraybackslash}X}
\newcommand{\headrow}{\rowcolor{snavy}}
\newcommand{\thc}[1]{\textcolor{white}{\bfseries #1}}
\newcommand{\ourrow}{\rowcolor{stint}}

% 图形
\usepackage{graphicx}
\usepackage{pgfplots}
\pgfplotsset{compat=1.18,tick label style={font=\numfont\small},label style={font=\small},/pgf/number format/assume math mode=true}
\tikzset{
  node/.style  = {draw=snavy, line width=0.8pt, fill=white, rounded corners=2pt, align=flush center,
                  inner xsep=7pt, inner ysep=5pt, font=\small},
  fillnode/.style = {fill=snavy, text=white, rounded corners=2pt, align=flush center,
                  inner xsep=7pt, inner ysep=5pt, font=\small\bfseries},
  tintnode/.style = {fill=stint, draw=none, rounded corners=2pt, align=flush center,
                  inner xsep=7pt, inner ysep=5pt, font=\small},
  key/.style   = {node, draw=sorange, line width=1.1pt},
  flow/.style  = {-{Stealth[length=2.6mm]}, line width=0.9pt, draw=snavy},
  note/.style  = {font=\small, text=sgrey, align=flush center},
}

% 封面：深蓝满版
\setbeamertemplate{title page}{%
  \begin{tikzpicture}[remember picture,overlay]
    \fill[snavy] (current page.south west) rectangle (current page.north east);
    \fill[sorange] ([xshift=1.2cm,yshift=-2.3cm]current page.north west) rectangle ++(1.6cm,0.12cm);
    \node[anchor=north west,text=white,text width=0.78\paperwidth,align=flush left]
      at ([xshift=1.2cm,yshift=-2.7cm]current page.north west)
      {{\titlecjk\fontsize{25pt}{33pt}\selectfont\inserttitle}\par\vspace{0.45cm}
       {\large\color{white!80!snavy}\insertsubtitle}};
    \node[anchor=south west,text=white,text width=0.6\paperwidth,align=flush left]
      at ([xshift=1.2cm,yshift=1.1cm]current page.south west)
      {{\large\bfseries\insertauthor}\par\vspace{0.12cm}{\small\color{white!75!snavy}\trackname\quad\insertdate}};
  \end{tikzpicture}}

% 章节页：深蓝满版，左侧大号编号与本章要回答的问题，右侧故事进度
\providecommand{\ChapterList}[1]{}
\newcommand{\ChapterFrame}[4]{%
  \section[#1\enspace #2]{#3}\gdef\chaplabel{#1\enspace #2}
  {\setbeamertemplate{footline}{}
  \begin{frame}[plain,noframenumbering]
    \begin{tikzpicture}[remember picture,overlay]
      \fill[snavy] (current page.south west) rectangle (current page.north east);
      \node[anchor=north west,text=white,text width=0.46\paperwidth,align=flush left]
        at ([xshift=1.2cm,yshift=-2.0cm]current page.north west)
        {{\numfont\color{sorange}\fontsize{48pt}{50pt}\selectfont #1}\par\vspace{0.35cm}
         {\titlecjk\fontsize{20pt}{26pt}\selectfont #3}\par\vspace{0.4cm}
         {\normalsize\color{white!85!snavy}#4}};
      \node[anchor=north west,text=white,text width=0.36\paperwidth,align=flush left]
        at ([xshift=0.58\paperwidth,yshift=-2.1cm]current page.north west)
        {\ChapterList{#1}};
    \end{tikzpicture}
  \end{frame}}}
\newcommand{\chapitem}[3]{\ifnum#1=#2{\color{white}\bfseries #2\enspace #3}\else{\color{white!55!snavy}#2\enspace #3}\fi\par\vspace{0.3cm}}

\hypersetup{colorlinks=true, linkcolor=snavy, urlcolor=snavy}

%   \title{…} \subtitle{…} \author{fx战队} \date{2026 年 11 月}
%   \TrackName{因子赛道}
%
% 页面结构：结论标题（浅显） + 主体（图 / 示意图 / 表 / 要点块） + 底部结论条 \takeaway{…}
% 组件：\ChapterFrame、\keynum、\point、panel 环境、\takeaway、\hl

\usepackage[fontset=none,scheme=plain]{ctex}
% 字体：微软雅黑（取自本机 Office 安装包，不复制进仓库），大号数字用 DIN
\def\yhpath{/Applications/Microsoft PowerPoint.app/Contents/Resources/DFonts/}
\setCJKmainfont{msyh.ttc}[Path=\yhpath,BoldFont=msyhbd.ttc]
\setCJKsansfont{msyh.ttc}[Path=\yhpath,BoldFont=msyhbd.ttc]
\setCJKmonofont{msyh.ttc}[Path=\yhpath]
\newCJKfontfamily\titlecjk{msyhbd.ttc}[Path=\yhpath]
\setsansfont{msyh.ttc}[Path=\yhpath,BoldFont=msyhbd.ttc]
\setmainfont{msyh.ttc}[Path=\yhpath,BoldFont=msyhbd.ttc]
\newfontfamily\titlelatin{msyhbd.ttc}[Path=\yhpath]
\newfontfamily\numfont{DIN Alternate}
\usefonttheme{professionalfonts}

\usepackage{newunicodechar}
\newunicodechar{→}{\ensuremath{\rightarrow}}
\newunicodechar{↑}{\ensuremath{\uparrow}}
\newunicodechar{↓}{\ensuremath{\downarrow}}
\newunicodechar{≈}{\ensuremath{\approx}}
\newunicodechar{≤}{\ensuremath{\leq}}
\newunicodechar{≥}{\ensuremath{\geq}}
\newunicodechar{−}{\ensuremath{-}}
\newunicodechar{×}{\ensuremath{\times}}
\newunicodechar{÷}{\ensuremath{\div}}
\newunicodechar{±}{\ensuremath{\pm}}
\newunicodechar{β}{\ensuremath{\beta}}
\newunicodechar{λ}{\ensuremath{\lambda}}

% 配色：深蓝主色、橙色强调、浅蓝底色
\definecolor{snavy}{HTML}{0F2E5A}
\definecolor{sblue}{HTML}{2F6DB5}
\definecolor{stint}{HTML}{E8EFF8}
\definecolor{sorange}{HTML}{E8772E}
\definecolor{sgrey}{HTML}{5E6672}
\definecolor{sline}{HTML}{C9D3E0}
\definecolor{sink}{HTML}{1E2329}

\setbeamercolor{normal text}{fg=sink,bg=white}
\setbeamercolor{structure}{fg=snavy}
\setbeamercolor{frametitle}{fg=snavy}
\setbeamercolor{alerted text}{fg=sorange}
\setbeamertemplate{navigation symbols}{}
\setbeamertemplate{itemize items}{\textcolor{sblue}{\rule[0.25ex]{0.42em}{0.42em}}}
\setlength{\leftmargini}{1.1em}
\linespread{1.18}
\setbeamersize{text margin left=1.0cm, text margin right=1.0cm}
% 内容顶端对齐，结论条用 \vfill 压到页底
\makeatletter\beamer@centeredfalse\makeatother

\providecommand{\trackname}{}
\gdef\chaplabel{}
\newcommand{\TrackName}[1]{\renewcommand{\trackname}{#1}}

\usepackage{tikz}
\usetikzlibrary{arrows.meta,positioning,calc,fit,decorations.pathreplacing,shapes.geometric}

% 顶部深蓝细带 + 标题左侧橙色色块 + 右上章节标签
\setbeamertemplate{frametitle}{%
  \begin{tikzpicture}[remember picture,overlay]
    \fill[snavy] (current page.north west) rectangle ([yshift=-0.12cm]current page.north east);
  \end{tikzpicture}%
  \vspace*{0.36cm}%
  \begin{minipage}[c]{0.80\textwidth}%
    {\color{sorange}\rule[-0.05cm]{0.13cm}{0.52cm}}\hspace{0.22cm}%
    {\titlelatin\titlecjk\color{snavy}\fontsize{16pt}{20pt}\selectfont\insertframetitle}%
  \end{minipage}\hfill
  \ifx\chaplabel\empty\else
  \begin{tikzpicture}[baseline=-0.6ex]
    \node[fill=snavy,text=white,rounded corners=1.5pt,inner xsep=6pt,inner ysep=3pt,font=\footnotesize]{\chaplabel};
  \end{tikzpicture}\fi\par
  \vspace{0.25cm}}

% 页脚：只有页码
\setbeamertemplate{footline}{\hfill{\numfont\color{sgrey}\small\insertframenumber}\hspace*{0.6cm}\vspace*{0.3cm}}

\usepackage{appendixnumberbeamer}

% ---------- 组件 ----------
\newcommand{\hl}[1]{\textcolor{sorange}{\bfseries #1}}
% 大号关键数字：\keynum{0.290}{换手率（原 0.518）}
\newcommand{\keynum}[2]{%
  \begin{minipage}[t]{\linewidth}\raggedright
    {\numfont\color{sorange}\fontsize{26pt}{28pt}\selectfont #1}\par\vspace{0.06cm}
    {\color{sink}\small #2}\par
  \end{minipage}}
% 编号要点：\point{1}{加粗短句}{一两句解释}
\newcommand{\point}[3]{%
  \begin{tikzpicture}[baseline=(n.base)]
    \node[circle,fill=snavy,text=white,inner sep=1.6pt,minimum size=0.48cm,font=\numfont\small](n){#1};
  \end{tikzpicture}\hspace{0.18cm}%
  \begin{minipage}[t]{\dimexpr\linewidth-0.75cm\relax}
    {\bfseries\color{snavy}#2}\par\vspace{0.03cm}{\small\color{sink}#3}
  \end{minipage}\par\vspace{0.16cm}}
% 浅蓝底面板：\begin{panel}{标题} … \end{panel}
\newenvironment{panel}[1]{%
  \par\noindent\begin{tikzpicture}\node[fill=stint,rounded corners=2pt,inner sep=8pt,text width=\dimexpr\linewidth-16pt\relax,align=flush left]\bgroup
  {\bfseries\color{snavy}#1}\par\vspace{0.12cm}\small}
  {\egroup;\end{tikzpicture}\par}
% 底部结论条：一句话说明这一页证明了什么（固定在页底，正文不要压到这一区域）
\newcommand{\takeaway}[1]{%
  \begin{tikzpicture}[remember picture,overlay]
    \node[fill=snavy,text=white,inner xsep=10pt,inner ysep=6pt,anchor=south,
          text width=\dimexpr\paperwidth-2cm-20pt\relax,align=flush left,font=\small]
      at ([yshift=0.75cm]current page.south) {#1};
  \end{tikzpicture}}

% 表格
\usepackage{booktabs,tabularx,array,colortbl}
\renewcommand{\arraystretch}{1.32}
\newcolumntype{L}{>{\raggedright\arraybackslash}X}
\newcolumntype{C}{>{\centering\arraybackslash}X}
\newcommand{\headrow}{\rowcolor{snavy}}
\newcommand{\thc}[1]{\textcolor{white}{\bfseries #1}}
\newcommand{\ourrow}{\rowcolor{stint}}

% 图形
\usepackage{graphicx}
\usepackage{pgfplots}
\pgfplotsset{compat=1.18,tick label style={font=\numfont\small},label style={font=\small},/pgf/number format/assume math mode=true}
\tikzset{
  node/.style  = {draw=snavy, line width=0.8pt, fill=white, rounded corners=2pt, align=flush center,
                  inner xsep=7pt, inner ysep=5pt, font=\small},
  fillnode/.style = {fill=snavy, text=white, rounded corners=2pt, align=flush center,
                  inner xsep=7pt, inner ysep=5pt, font=\small\bfseries},
  tintnode/.style = {fill=stint, draw=none, rounded corners=2pt, align=flush center,
                  inner xsep=7pt, inner ysep=5pt, font=\small},
  key/.style   = {node, draw=sorange, line width=1.1pt},
  flow/.style  = {-{Stealth[length=2.6mm]}, line width=0.9pt, draw=snavy},
  note/.style  = {font=\small, text=sgrey, align=flush center},
}

% 封面：深蓝满版
\setbeamertemplate{title page}{%
  \begin{tikzpicture}[remember picture,overlay]
    \fill[snavy] (current page.south west) rectangle (current page.north east);
    \fill[sorange] ([xshift=1.2cm,yshift=-2.3cm]current page.north west) rectangle ++(1.6cm,0.12cm);
    \node[anchor=north west,text=white,text width=0.78\paperwidth,align=flush left]
      at ([xshift=1.2cm,yshift=-2.7cm]current page.north west)
      {{\titlecjk\fontsize{25pt}{33pt}\selectfont\inserttitle}\par\vspace{0.45cm}
       {\large\color{white!80!snavy}\insertsubtitle}};
    \node[anchor=south west,text=white,text width=0.6\paperwidth,align=flush left]
      at ([xshift=1.2cm,yshift=1.1cm]current page.south west)
      {{\large\bfseries\insertauthor}\par\vspace{0.12cm}{\small\color{white!75!snavy}\trackname\quad\insertdate}};
  \end{tikzpicture}}

% 章节页：深蓝满版，左侧大号编号与本章要回答的问题，右侧故事进度
\providecommand{\ChapterList}[1]{}
\newcommand{\ChapterFrame}[4]{%
  \section[#1\enspace #2]{#3}\gdef\chaplabel{#1\enspace #2}
  {\setbeamertemplate{footline}{}
  \begin{frame}[plain,noframenumbering]
    \begin{tikzpicture}[remember picture,overlay]
      \fill[snavy] (current page.south west) rectangle (current page.north east);
      \node[anchor=north west,text=white,text width=0.46\paperwidth,align=flush left]
        at ([xshift=1.2cm,yshift=-2.0cm]current page.north west)
        {{\numfont\color{sorange}\fontsize{48pt}{50pt}\selectfont #1}\par\vspace{0.35cm}
         {\titlecjk\fontsize{20pt}{26pt}\selectfont #3}\par\vspace{0.4cm}
         {\normalsize\color{white!85!snavy}#4}};
      \node[anchor=north west,text=white,text width=0.36\paperwidth,align=flush left]
        at ([xshift=0.58\paperwidth,yshift=-2.1cm]current page.north west)
        {\ChapterList{#1}};
    \end{tikzpicture}
  \end{frame}}}
\newcommand{\chapitem}[3]{\ifnum#1=#2{\color{white}\bfseries #2\enspace #3}\else{\color{white!55!snavy}#2\enspace #3}\fi\par\vspace{0.3cm}}

\hypersetup{colorlinks=true, linkcolor=snavy, urlcolor=snavy}

%   \title{…} \subtitle{…} \author{fx战队} \date{2026 年 11 月}
%   \TrackName{因子赛道}
%
% 页面结构：结论标题（浅显） + 主体（图 / 示意图 / 表 / 要点块） + 底部结论条 \takeaway{…}
% 组件：\ChapterFrame、\keynum、\point、panel 环境、\takeaway、\hl

\usepackage[fontset=none,scheme=plain]{ctex}
% 字体：微软雅黑（取自本机 Office 安装包，不复制进仓库），大号数字用 DIN
\def\yhpath{/Applications/Microsoft PowerPoint.app/Contents/Resources/DFonts/}
\setCJKmainfont{msyh.ttc}[Path=\yhpath,BoldFont=msyhbd.ttc]
\setCJKsansfont{msyh.ttc}[Path=\yhpath,BoldFont=msyhbd.ttc]
\setCJKmonofont{msyh.ttc}[Path=\yhpath]
\newCJKfontfamily\titlecjk{msyhbd.ttc}[Path=\yhpath]
\setsansfont{msyh.ttc}[Path=\yhpath,BoldFont=msyhbd.ttc]
\setmainfont{msyh.ttc}[Path=\yhpath,BoldFont=msyhbd.ttc]
\newfontfamily\titlelatin{msyhbd.ttc}[Path=\yhpath]
\newfontfamily\numfont{DIN Alternate}
\usefonttheme{professionalfonts}

\usepackage{newunicodechar}
\newunicodechar{→}{\ensuremath{\rightarrow}}
\newunicodechar{↑}{\ensuremath{\uparrow}}
\newunicodechar{↓}{\ensuremath{\downarrow}}
\newunicodechar{≈}{\ensuremath{\approx}}
\newunicodechar{≤}{\ensuremath{\leq}}
\newunicodechar{≥}{\ensuremath{\geq}}
\newunicodechar{−}{\ensuremath{-}}
\newunicodechar{×}{\ensuremath{\times}}
\newunicodechar{÷}{\ensuremath{\div}}
\newunicodechar{±}{\ensuremath{\pm}}
\newunicodechar{β}{\ensuremath{\beta}}
\newunicodechar{λ}{\ensuremath{\lambda}}

% 配色：深蓝主色、橙色强调、浅蓝底色
\definecolor{snavy}{HTML}{0F2E5A}
\definecolor{sblue}{HTML}{2F6DB5}
\definecolor{stint}{HTML}{E8EFF8}
\definecolor{sorange}{HTML}{E8772E}
\definecolor{sgrey}{HTML}{5E6672}
\definecolor{sline}{HTML}{C9D3E0}
\definecolor{sink}{HTML}{1E2329}

\setbeamercolor{normal text}{fg=sink,bg=white}
\setbeamercolor{structure}{fg=snavy}
\setbeamercolor{frametitle}{fg=snavy}
\setbeamercolor{alerted text}{fg=sorange}
\setbeamertemplate{navigation symbols}{}
\setbeamertemplate{itemize items}{\textcolor{sblue}{\rule[0.25ex]{0.42em}{0.42em}}}
\setlength{\leftmargini}{1.1em}
\linespread{1.18}
\setbeamersize{text margin left=1.0cm, text margin right=1.0cm}
% 内容顶端对齐，结论条用 \vfill 压到页底
\makeatletter\beamer@centeredfalse\makeatother

\providecommand{\trackname}{}
\gdef\chaplabel{}
\newcommand{\TrackName}[1]{\renewcommand{\trackname}{#1}}

\usepackage{tikz}
\usetikzlibrary{arrows.meta,positioning,calc,fit,decorations.pathreplacing,shapes.geometric}

% 顶部深蓝细带 + 标题左侧橙色色块 + 右上章节标签
\setbeamertemplate{frametitle}{%
  \begin{tikzpicture}[remember picture,overlay]
    \fill[snavy] (current page.north west) rectangle ([yshift=-0.12cm]current page.north east);
  \end{tikzpicture}%
  \vspace*{0.36cm}%
  \begin{minipage}[c]{0.80\textwidth}%
    {\color{sorange}\rule[-0.05cm]{0.13cm}{0.52cm}}\hspace{0.22cm}%
    {\titlelatin\titlecjk\color{snavy}\fontsize{16pt}{20pt}\selectfont\insertframetitle}%
  \end{minipage}\hfill
  \ifx\chaplabel\empty\else
  \begin{tikzpicture}[baseline=-0.6ex]
    \node[fill=snavy,text=white,rounded corners=1.5pt,inner xsep=6pt,inner ysep=3pt,font=\footnotesize]{\chaplabel};
  \end{tikzpicture}\fi\par
  \vspace{0.25cm}}

% 页脚：只有页码
\setbeamertemplate{footline}{\hfill{\numfont\color{sgrey}\small\insertframenumber}\hspace*{0.6cm}\vspace*{0.3cm}}

\usepackage{appendixnumberbeamer}

% ---------- 组件 ----------
\newcommand{\hl}[1]{\textcolor{sorange}{\bfseries #1}}
% 大号关键数字：\keynum{0.290}{换手率（原 0.518）}
\newcommand{\keynum}[2]{%
  \begin{minipage}[t]{\linewidth}\raggedright
    {\numfont\color{sorange}\fontsize{26pt}{28pt}\selectfont #1}\par\vspace{0.06cm}
    {\color{sink}\small #2}\par
  \end{minipage}}
% 编号要点：\point{1}{加粗短句}{一两句解释}
\newcommand{\point}[3]{%
  \begin{tikzpicture}[baseline=(n.base)]
    \node[circle,fill=snavy,text=white,inner sep=1.6pt,minimum size=0.48cm,font=\numfont\small](n){#1};
  \end{tikzpicture}\hspace{0.18cm}%
  \begin{minipage}[t]{\dimexpr\linewidth-0.75cm\relax}
    {\bfseries\color{snavy}#2}\par\vspace{0.03cm}{\small\color{sink}#3}
  \end{minipage}\par\vspace{0.16cm}}
% 浅蓝底面板：\begin{panel}{标题} … \end{panel}
\newenvironment{panel}[1]{%
  \par\noindent\begin{tikzpicture}\node[fill=stint,rounded corners=2pt,inner sep=8pt,text width=\dimexpr\linewidth-16pt\relax,align=flush left]\bgroup
  {\bfseries\color{snavy}#1}\par\vspace{0.12cm}\small}
  {\egroup;\end{tikzpicture}\par}
% 底部结论条：一句话说明这一页证明了什么（固定在页底，正文不要压到这一区域）
\newcommand{\takeaway}[1]{%
  \begin{tikzpicture}[remember picture,overlay]
    \node[fill=snavy,text=white,inner xsep=10pt,inner ysep=6pt,anchor=south,
          text width=\dimexpr\paperwidth-2cm-20pt\relax,align=flush left,font=\small]
      at ([yshift=0.75cm]current page.south) {#1};
  \end{tikzpicture}}

% 表格
\usepackage{booktabs,tabularx,array,colortbl}
\renewcommand{\arraystretch}{1.32}
\newcolumntype{L}{>{\raggedright\arraybackslash}X}
\newcolumntype{C}{>{\centering\arraybackslash}X}
\newcommand{\headrow}{\rowcolor{snavy}}
\newcommand{\thc}[1]{\textcolor{white}{\bfseries #1}}
\newcommand{\ourrow}{\rowcolor{stint}}

% 图形
\usepackage{graphicx}
\usepackage{pgfplots}
\pgfplotsset{compat=1.18,tick label style={font=\numfont\small},label style={font=\small},/pgf/number format/assume math mode=true}
\tikzset{
  node/.style  = {draw=snavy, line width=0.8pt, fill=white, rounded corners=2pt, align=flush center,
                  inner xsep=7pt, inner ysep=5pt, font=\small},
  fillnode/.style = {fill=snavy, text=white, rounded corners=2pt, align=flush center,
                  inner xsep=7pt, inner ysep=5pt, font=\small\bfseries},
  tintnode/.style = {fill=stint, draw=none, rounded corners=2pt, align=flush center,
                  inner xsep=7pt, inner ysep=5pt, font=\small},
  key/.style   = {node, draw=sorange, line width=1.1pt},
  flow/.style  = {-{Stealth[length=2.6mm]}, line width=0.9pt, draw=snavy},
  note/.style  = {font=\small, text=sgrey, align=flush center},
}

% 封面：深蓝满版
\setbeamertemplate{title page}{%
  \begin{tikzpicture}[remember picture,overlay]
    \fill[snavy] (current page.south west) rectangle (current page.north east);
    \fill[sorange] ([xshift=1.2cm,yshift=-2.3cm]current page.north west) rectangle ++(1.6cm,0.12cm);
    \node[anchor=north west,text=white,text width=0.78\paperwidth,align=flush left]
      at ([xshift=1.2cm,yshift=-2.7cm]current page.north west)
      {{\titlecjk\fontsize{25pt}{33pt}\selectfont\inserttitle}\par\vspace{0.45cm}
       {\large\color{white!80!snavy}\insertsubtitle}};
    \node[anchor=south west,text=white,text width=0.6\paperwidth,align=flush left]
      at ([xshift=1.2cm,yshift=1.1cm]current page.south west)
      {{\large\bfseries\insertauthor}\par\vspace{0.12cm}{\small\color{white!75!snavy}\trackname\quad\insertdate}};
  \end{tikzpicture}}

% 章节页：深蓝满版，左侧大号编号与本章要回答的问题，右侧故事进度
\providecommand{\ChapterList}[1]{}
\newcommand{\ChapterFrame}[4]{%
  \section[#1\enspace #2]{#3}\gdef\chaplabel{#1\enspace #2}
  {\setbeamertemplate{footline}{}
  \begin{frame}[plain,noframenumbering]
    \begin{tikzpicture}[remember picture,overlay]
      \fill[snavy] (current page.south west) rectangle (current page.north east);
      \node[anchor=north west,text=white,text width=0.46\paperwidth,align=flush left]
        at ([xshift=1.2cm,yshift=-2.0cm]current page.north west)
        {{\numfont\color{sorange}\fontsize{48pt}{50pt}\selectfont #1}\par\vspace{0.35cm}
         {\titlecjk\fontsize{20pt}{26pt}\selectfont #3}\par\vspace{0.4cm}
         {\normalsize\color{white!85!snavy}#4}};
      \node[anchor=north west,text=white,text width=0.36\paperwidth,align=flush left]
        at ([xshift=0.58\paperwidth,yshift=-2.1cm]current page.north west)
        {\ChapterList{#1}};
    \end{tikzpicture}
  \end{frame}}}
\newcommand{\chapitem}[3]{\ifnum#1=#2{\color{white}\bfseries #2\enspace #3}\else{\color{white!55!snavy}#2\enspace #3}\fi\par\vspace{0.3cm}}

\hypersetup{colorlinks=true, linkcolor=snavy, urlcolor=snavy}
